\uuid{jXw7}
\exo7id{1732}
\titre{exo7 1732}
\auteur{ridde}
\organisation{exo7}
\datecreate{1999-11-01}
\isIndication{false}
\isCorrection{false}
\chapitre{Fonction convexe}
\sousChapitre{Fonction convexe}
\cpge{mpsi-mp2i.convexite.a, mpsi-mp2i.convexite.b, mpsi-mp2i.derivabilite.c, pcsi.derivabilite.c, ptsi.derivabilite.c, mpsi-mp2i.analyse-asymptotique.b, pcsi.analyse-asymptotique.b, ptsi.analyse-asymptotique.b, bcpst1.limites-continuite.a, bcpst1.derivation.a, bcpst1.derivation.b, bcpst1.developpements-limites.b}
\module{Optimisation}
\niveau{L3}
\difficulte{}

\contenu{
\texte{
Soit $f \in C^2 (\Rr)$ convexe.
}
\begin{enumerate}
    \item \question{Montrer que $f'$ admet une limite dans $\bar{\Rr}$ en $ + \infty$.}
    \item \question{En d\'eduire que $\frac{f (x)}x$ admet une limite en $ + \infty$
 (on pourra utiliser des $\epsilon$ et une formule de Taylor \`a l'ordre 1).}
\end{enumerate}
}
